\documentclass[11pt]{article}
\usepackage[a4paper, left=1.5cm,right=1.5cm,top=2cm,bottom=2cm]{geometry}
\usepackage{graphicx} % Required for inserting images
\usepackage{amsmath}  % Math environments and symbols
\usepackage{amssymb}  % \mathbb and extra symbols
\usepackage{booktabs} % Nicer tables (\toprule, \midrule, \bottomrule)

\title{Compositional Latent Space}
\author{Suresh Murugaiyan}
\date{July 2026}

\begin{document}

\maketitle

\section{Introduction}

We train a \emph{compositional autoencoder} (Path~A of the compositional
operator-network framework) on the FlowBench 2D lid-driven cavity dataset:
steady Navier--Stokes flow with an object placed inside the cavity. The
autoencoder compresses each steady flow field into a small latent vector
$\boldsymbol{z}$ and reconstructs the field from it --- a nonlinear analogue
of POD. The key structural choice is that $\boldsymbol{z}$ is not one opaque
vector; it is partitioned by architecture into three named blocks, each
intended to carry one physical factor:
%
\begin{equation}
  \boldsymbol{z} \;=\;
  [\, \underbrace{\boldsymbol{z}_{\mu}}_{\text{regime }(4)}
  \,\|\,
  \underbrace{\boldsymbol{z}_{g}}_{\text{geometry }(32)}
  \,\|\,
  \underbrace{\boldsymbol{z}_{\xi}}_{\text{residual }(16)} \,]
  \;\in\; \mathbb{R}^{52}.
\end{equation}
%
The hypothesis under test: after training, $\boldsymbol{z}_{\mu}$ should
contain the Reynolds number, $\boldsymbol{z}_{g}$ should contain the object
shape, and neither should leak into the other. Because the flow is steady,
there is no dynamics block $\boldsymbol{z}_{\eta}$ and no latent propagator
$\Phi$ in this first version; those enter later with time-dependent cases.

\section{Dataset}

Each FlowBench sample is a steady lid-driven cavity solution at
$512 \times 512$ resolution. The data loader
(\texttt{data/dataset.py}) unpacks the raw \texttt{.npz} arrays into:
%
\begin{itemize}
  \item \textbf{fields} --- the $(u, v, p)$ solution, resized
        $512^2 \to 256^2$ to keep training cheap;
  \item \textbf{sdf} --- the signed distance field of the object
        (the geometry label);
  \item \textbf{mask} --- binary fluid mask ($1$ in fluid, $0$ in the solid);
  \item \textbf{log\_re} --- the regime label. Re is stored as a constant
        channel, so a single pixel is extracted, transformed to
        $\log_{10}\mathrm{Re}$ (Re spans orders of magnitude), and
        standardized with \emph{train-set} statistics; the test set reuses
        those statistics so both splits share one scale.
\end{itemize}
%
The loader also computes three cheap per-sample geometry summaries from the
mask (solid area fraction, centroid $x$, centroid $y$). These are never used
during training; they serve only as ground-truth targets for the linear-probe
diagnostic of Section~\ref{sec:diagnostics}.

\section{Network architecture}

Four modules (\texttt{models/compositional/networks.py}):
%
\begin{itemize}
  \item \textbf{FieldEncoder} $\mathcal{E}$ --- a CNN that halves the spatial
        resolution repeatedly ($256 \to 8$), pools, and terminates in
        \emph{three separate linear heads} producing
        $\boldsymbol{z}_{\mu}, \boldsymbol{z}_{g}, \boldsymbol{z}_{\xi}$.
        The blocks are architecturally separate from the start.
  \item \textbf{FieldDecoder} $\mathcal{D}$ --- the mirror image: the
        concatenated $\boldsymbol{z}$ is reshaped to an $8\times 8$ feature
        map and upsampled back to a $256^2$ $(u, v, p)$ field.
  \item \textbf{RegimeHead} --- a small MLP that reads
        $\boldsymbol{z}_{\mu}$ \emph{only} and must predict
        $\log \mathrm{Re}$.
  \item \textbf{SDFHead} --- a small decoder that reads $\boldsymbol{z}_{g}$
        \emph{only} and must redraw the object's SDF at $64 \times 64$.
\end{itemize}
%
The two heads are the supervision mechanism: since the regime head sees only
$\boldsymbol{z}_{\mu}$, the encoder is forced to route Reynolds-number
information into that block, and likewise shape information into
$\boldsymbol{z}_{g}$.

\section{Training objective}

Each optimization step combines four losses
(\texttt{models/compositional/compositional\_ae.py}), weighted by
$\lambda$-coefficients set in the configuration file:
%
\begin{enumerate}
  \item \textbf{Masked reconstruction (L1).} MSE between decoded and true
        fields, evaluated over fluid points only: predictions inside the
        solid are zeroed and the error is normalized by the fluid node
        count, so samples with large objects are not unfairly weighted.
  \item \textbf{Regime supervision.} MSE of
        $\mathrm{RegimeHead}(\boldsymbol{z}_{\mu})$ against the standardized
        $\log_{10}\mathrm{Re}$. Pulls Re into $\boldsymbol{z}_{\mu}$.
  \item \textbf{Geometry supervision.} MSE of
        $\mathrm{SDFHead}(\boldsymbol{z}_{g})$ against the downsampled true
        SDF. Pulls shape into $\boldsymbol{z}_{g}$.
  \item \textbf{Cross-block decorrelation (L6).} The mean absolute Pearson
        correlation, computed across the batch, between every coordinate of
        one block and every coordinate of another, penalized for all three
        block pairs. Losses~2--3 pull information \emph{into} the correct
        block; this term pushes redundant copies \emph{out of} the wrong
        blocks.
\end{enumerate}
%
The residual block $\boldsymbol{z}_{\xi}$ receives no supervision at all: it
is free capacity that absorbs whatever the reconstruction needs beyond
$(\mathrm{Re}, \text{shape})$, which is exactly the role of the residual
block in the framework.

Training is driven by \texttt{main.py}: it loads the YAML configuration,
seeds all libraries, builds the train/test datasets (passing the train
Re-statistics to the test set), and trains with PyTorch Lightning,
checkpointing on the best validation reconstruction error.

\section{Linear-probe diagnostics}
\label{sec:diagnostics}

The diagnostic (\texttt{diagnostics/probes.py}) is the actual scientific
test. After training, the entire test set is encoded, and for each
(block, target) pair a ridge regression is fitted with 5-fold
cross-validation; the resulting $R^2$ values are tabulated. For a
compositional latent we expect the block-diagonal pattern of
Table~\ref{tab:probes}.

\begin{table}[htbp]
  \centering
  \caption{Expected linear-probe $R^2$ pattern for a compositional latent
  space. High $R^2$ on the matching block and low $R^2$ elsewhere indicates
  a successful factorization; a high off-diagonal entry indicates leakage.}
  \label{tab:probes}
  \begin{tabular}{lccc}
    \toprule
    Target & $\boldsymbol{z}_{\mu}$ & $\boldsymbol{z}_{g}$ & $\boldsymbol{z}_{\xi}$ \\
    \midrule
    $\log \mathrm{Re}$    & \textbf{high} & low           & low \\
    Area fraction         & low           & \textbf{high} & low \\
    Centroid $x$          & low           & \textbf{high} & low \\
    Centroid $y$          & low           & \textbf{high} & low \\
    \bottomrule
  \end{tabular}
\end{table}

If the trained latent reproduces this pattern, the factorization worked:
Re is linearly readable from the regime block and \emph{not} from the
geometry block, and vice versa. A high $R^2$ in an off-diagonal cell means
information leaked --- exactly what the planned loss ablations are designed
to study (e.g., does removing the decorrelation penalty cause leakage?).

\section{Next steps}

Deliberately excluded from this first version and planned as increments once
the baseline trains and the probe table looks sane:
%
\begin{itemize}
  \item swap-consistency loss (L12): decode $\boldsymbol{z}_{\mu}$ from one
        sample with $\boldsymbol{z}_{g}$ from another and compare against
        the true cross-combination;
  \item group-structured minibatches for iVAE-style identifiability;
  \item concept-vector arithmetic diagnostic;
  \item INR-style implicit decoder; HSIC (nonlinear) decorrelation;
  \item dynamics block $\boldsymbol{z}_{\eta}$ and latent propagator $\Phi$
        for time-dependent cases.
\end{itemize}

\end{document}
